% 2026 全国 2 卷 15：时间（天）的频率分布直方图
% 八组，组距均为 10（[345,355),…,[415,425)），频率/组距依次为
% 0.005、0.010、0.020、0.025、0.015、0.015、0.005、0.005（合计 ×10 = 1）。
% 第 5、6 组等高，第 7、8 组等高，故轮廓在这两处连成一条顶边（不画中间竖线）。
% 横轴上 O 与 345 之间画波浪断轴符号；纵轴各刻度处的水平虚线向右贯穿全图。
% 绘图坐标以「组序」为单位：横坐标 t 对应实际值 345+10t，原点 O 置于 t=-1.1。
\begin{tikzpicture}[font=\small, x=0.90cm, y=200cm]
  % ---- 水平虚线网格 ----
  \foreach \y in {0.005,0.010,0.015,0.020,0.025}
    { \draw[dashed] (-1.1,\y) -- (8.35,\y); }
  % ---- 直方图轮廓（等高相邻组不画中间竖线）----
  \draw (0,0) -- (0,0.005) -- (1,0.005) -- (1,0.010) -- (2,0.010) -- (2,0.020)
        -- (3,0.020) -- (3,0.025) -- (4,0.025) -- (4,0.015) -- (6,0.015)
        -- (6,0.005) -- (8,0.005) -- (8,0);
  % ---- 坐标轴（横轴在断轴处留缺口）----
  \draw (-1.1,0) -- (-0.85,0);
  \draw[-stealth] (-0.39,0) -- (8.6,0);
  \draw[-stealth] (-1.1,0) -- (-1.1,0.0272);
  % ---- 波浪断轴符号 ----
  \begin{scope}[shift={(-0.62,0)}]
    \begin{scope}[x=1cm,y=1cm]
      \draw (-0.22,0) .. controls (-0.13,0.13) and (-0.05,-0.13) .. (0,0)
                      .. controls (0.05,0.13)  and (0.13,-0.13) .. (0.22,0);
    \end{scope}
  \end{scope}
  % ---- 刻度 ----
  \foreach \t/\lab in {0/345,1/355,2/365,3/375,4/385,5/395,6/405,7/415,8/425}
    { \draw (\t,0) -- (\t,-0.0006); \node[below, inner sep=2pt] at (\t,-0.0004) {\footnotesize $\lab$}; }
  \foreach \y/\lab in {0.005/0.005,0.010/0.01,0.015/0.015,0.020/0.02,0.025/0.025}
    { \draw (-1.1,\y) -- (-1.22,\y); \node[left, inner sep=2pt] at (-1.22,\y) {\footnotesize $\lab$}; }
  \node[below left, inner sep=1.5pt] at (-1.1,0) {$O$};
  % ---- 轴名 ----
  \node[right, inner sep=3pt] at (8.6,0) {时间/天};
  \node[above, inner sep=2.5pt] at (-1.1,0.0272) {$\displaystyle\frac{\mbox{频率}}{\mbox{组距}}$};
\end{tikzpicture}
